\documentclass[10pt,a4paper]{amsart}
\usepackage[margin=19mm]{geometry}
\usepackage{amsmath,amssymb,mathtools}
\usepackage[scr=rsfs,cal=euler]{mathalfa}
\usepackage{xfrac}
\usepackage[unicode,hidelinks,bookmarks=false]{hyperref}
\newcommand{\ie}{i.e.\ }
\newcommand{\cf}{cf.\ }
\newcommand{\resp}{respectively\ }
\newcommand{\Subsection}{\subsection}
\providecommand{\nobreakdash}{\nobreak\hbox{-}}
\setlength{\emergencystretch}{2em}
\multlinegap0pt
\usepackage{bm}
\usepackage{bbm}
\usepackage{dsfont}
\usepackage{xspace}
\usepackage{cancel}
\usepackage{mathtools}
\usepackage{xcolor}
\usepackage{amsthm}
\makeatletter
\@ifundefined{theorem}{\newtheorem{theorem}{Theorem}}{}
\@ifundefined{lemma}{\newtheorem{lemma}[theorem]{Lemma}}{}
\makeatother
\newcommand{\thecolor}{black} 
\title[Structural description of (bull, house)-free graphs]{Structural description of (bull, house)-free graphs}
\author{Manoj Belavadi, Chinh T. Hoang}
\date{Source: arXiv:2604.27594v3 (CC BY 4.0)}
\begin{document}
\maketitle

\noindent\emph{Source and licence.} Structural description of (bull, house)-free graphs. Version arXiv:2604.27594v3 (CC BY 4.0), distributed under the Creative Commons Attribution 4.0 International licence, as recorded in the arXiv licence metadata for this exact version and shown on the abstract page https://arxiv.org/abs/2604.27594v3. Full rights notice: licenses/arxiv-2604.27594-CC-BY-4.0.txt.\par\smallskip

\noindent\textit{Editorial scope.} Counted unit: the Lemma whose source label is \texttt{lem:hole} in the article (source0.tex lines 248--250, source label \texttt{lem:hole}), delivered together with the article's own complete original proof of it (source0.tex lines 251--285, one \texttt{proof} environment of 35 lines). No mathematics is written, repaired or re-derived: statement and proof are the article's own text line for line. Required context retained uncounted: lem:split-vertex (lines 230--232). Every definition, display and label of those lines is retained unchanged. Omitted from this package and not redistributed: 1--229, 233--247, 286--703. Nothing inside a retained excerpt is omitted or altered: every retained body is delivered exactly as the article's lines are written, so no omission receipt of any kind accompanies this package. Citations: the retained excerpts carry no citation command at all, so no bibliography entry of the article is transcribed and no reference is redistributed. Reference adaptation: none was needed and none was made; every reference inside the delivered unit resolves inside it, so no unresolved reference or citation is produced. Printed slips kept literal: no printed word, symbol, hypothesis or number of the counted statement or of its proof was altered. Layout and class: the article's own class is \texttt{article} with packages including geometry, amsmath, amsthm, amssymb, tikz, enumitem, comment, hyperref, xcolor; the wrapper is the lane's pinned \texttt{amsart} editorial wrapper at 10pt on A4, which restates the theorem-like environments the retained text needs and adds the article's own macro definitions that the retained text needs (\texttt{\textbackslash thecolor}), with extra wrapper packages (bm, bbm, dsfont, xspace, cancel, mathtools, xcolor). The delivered numbering is the unit's own. Assets: no retained span contains a figure environment, a graphics-inclusion command, a PDF-page inclusion, a table or a TikZ picture; no image file is copied into this package, the package declares no asset dependency and no image file is redistributed with it. Licence: version arXiv:2604.27594v3 of this article is distributed under the Creative Commons Attribution 4.0 International licence, as recorded in the arXiv licence metadata for this exact version and shown on the abstract page https://arxiv.org/abs/2604.27594v3. A full rights notice is delivered at licenses/arxiv-2604.27594-CC-BY-4.0.txt. Characters: every retained excerpt is pure ASCII, so the delivered engine, XeLaTeX with its default Latin Modern text font, reports no missing character.
\begin{lemma}\label{lem:split-vertex}
    Let $G$ be a connected (bull, house)-free graph. Suppose $v$ is a vertex of $G$, such that $G[N(v)]$ has two components, then $G$ is triangle-free or $G$ has a  homogeneous set \textcolor{\thecolor}{$H$ such that $G[H]$ is connected}.
\end{lemma}

\begin{lemma}\label{lem:hole}
    Let $G$ be a connected (bull, house)-free graph. Suppose $G$ contains a hole with at least five vertices, then $G$ is triangle-free, or $G$ contains a  homogeneous set $H$ such that $G[H]$ is connected.
\end{lemma}

\begin{proof}
    Assume $G$ contains a hole induced by $X = \{x_1, x_2, \ldots, x_{\ell}\}$ such that $x_{i}x_{i+1}\in E(G)$ where $i$ is taken modulo $ \ell \ge 5$. Let $U$ be the set of vertices
    adjacent to every vertex in $X$. Suppose $U$ is nonempty. We first prove that $U$ is complete to $N(X)\setminus U$. Let $y\in N(X)\setminus U$, then there exists $j$ modulo $l$ such that $y \sim x_j$  but  $y \not\sim x_{j-1}$. Suppose $a\in U$ is not adjacent to $y$, then $\{y, x_{j}, x_{j-1}, a, x_{j+2}\}$ induces either a bull or a house, a contradiction. So $U$ is complete to $N(X)\setminus U$.
    
    Write $N_0 = X$, $N_1 = N(X)\setminus U$. Define, for $i \geq 2$,  $N_i = \{ v \in V(G) \setminus (U \cup (\cup_{k=0)}^{k=i-1} N_k)) \; | \; v$ has a neighbor in $N_{i-1} \}$.

    We will prove that  
\begin{equation}\label{eq:level}
	\text{\parbox{.85\textwidth}{$U$ is complete to $\cup_{k=0}^{i} N_k$ for any index $i \geq 0$}}
\end{equation}
    We prove (\ref{eq:level}) by induction on $i$. We know (\ref{eq:level}) holds for $i = 0,1$. Consider the case $i=2$. Suppose that there is a vertex $a \in U$ that is not adjacent to a vertex $z \in N_2$. Consider a neighbor $y \in N_1$ of $z$. We know $a \sim y$. There is an index $j$ such that  $ y \sim x_j$ and  $ y \not\sim x_{j-1}$. Vertex $y$ is adjacent to $x_{j-2}$, for otherwise $\{ x_j, y, a, x_{j-2}, z \}$ induces a bull. 

%textcolor{blue}{   ******  
 %we have not defined the `order' of the bull! But if we do, I would prefer to list the triangle, then the two pendant vertices. This helps the reader to see the bull. I would do the same for the house. 
% }

Vertex $y$ is not adjacent to $x_{j+1}$, for otherwise $\{x_{j+1}, y, x_j, x_{j-1}, x_{j-2}  \}$ induces a house. But now $\{ x_{j-2}, y, a, x_{j+1}, z \}$ induces a bull. So $U$ is complete to $N_2$. Now, by the induction hypothesis, we may assume $U$ is complete to $\cup_{k=0}^{k=i-1} N_k$. For $i \geq 3$, suppose there is a vertex $a \in U$ that is not adjacent to a vertex $z \in N_i$. Consider vertices $b_{i-1} \in N_{i-1}$, $b_{i-2} \in N_{i-2}$ such that $b_{i-1} \sim b_{i-2}, z \sim b_{i-1}$, $z \not\sim b_{i-2}$. Consider a vertex $r \in X$ that is not adjacent to $b_{i-2}$. Now $\{b_{i-2}, b_{i-1}, a, r, z \}$ induces a bull. So (\ref{eq:level}) holds. 
    
    Let $R = V(G) \setminus U \cup (\cup_{k=0}^{i} N_k)$. Then $R$ is anti-complete to  $\cup_{k=0}^{i} N_k$. Thus, $\cup_{k=0}^{i} N_k$ is a connected homogeneous set. So we may assume $U$ is empty. 

    \begin{equation}\label{eq:three-consecutive}
	\text{\parbox{.85\textwidth}{Let $y$ be a vertex in $V(G) \setminus X$ with $y \sim x_i, y \sim x_{i+1}$. Then we have $N(y)\cap X = \{x_i,x_{i+1},x_{i+2}\}$ or $N(y) \cap X = \{x_{i-1}, x_i,x_{i+1}\}$.}}
    \end{equation}
   

 For simplicity, we will let $i=2$. Let $y$ be a vertex in $V(G) \setminus X$ with $y \sim x_2, y \sim x_{3}$.  Suppose that (\ref{eq:three-consecutive}) does not hold for $y$.
 Since $\{y,x_2, x_3, x_{4}, x_1\}$ does not induce a bull, $y$ must be adjacent to $x_1$ or $x_4$. Assume $y \sim x_{1}$. Let $j$ be the smallest index such that $yx_j\notin E(G)$ (since $U = \emptyset$, $j$ exists) and let $k \geq 4 $ be the largest index such that $y x_k\in E(G)$ ($k$ exists since (\ref{eq:three-consecutive}) fails for $y$). If $j = 4$, then $\{ x_2, x_3, y, x_k, x_4 \}$ induces a bull or a house. So we have $j > 4$. But now $\{ x_{j-2}  , x_{j-1} , y, x_1, x_j\}$ induces a bull or a house. So (\ref{eq:three-consecutive}) holds. 
 
    %

Suppose $N(x_2)\cap N(x_3)$ is empty, then $G[N(x_2)]$ is disconnected, and the result follows from Lemma \ref{lem:split-vertex}. So assume $N(x_2)\cap N(x_3)$ is non-empty. Note that, from (\ref{eq:three-consecutive}), the set $N(x_2)\cap N(x_3)$ is contained in $A_2 \cup A_3$, where $A_2$ = $\{v\in V(G)| N(v)\cap X = \{x_1, x_2, x_3\} \text{ or } N(v)\cap X = \{x_1, x_3\} \}$ and $A_3$ = $\{v\in V(G)| N(v)\cap X = \{x_2, x_3, x_4\} \text{ or } N(v)\cap X = \{x_2, x_4\}\}$. Note that $x_2\in A_2$ and $x_3\in A_3$.
    
    
    We may assume there is a vertex $y\in A_2$ such that $yx_2\in E(G)$. Let $B$ induce the component in $G[A_2]$ that contains the edge $yx_2$. We claim that $B$ is a homogeneous set. Suppose not, there exist adjacent vertices $b_1$, $b_2$ in $B$ such that there exists a vertex $z\notin B$ adjacent to $b_2$ but not to $b_1$. If $z$ is not adjacent to $x_3$ (resp. $x_1$), then $\{  b_1, b_2, x_3, x_4, z\}$ (resp. $\{ b_1, b_2,x_1, x_\ell,  z\}$) induces either a bull or a house. Hence $z\in N(x_1)\cap N(x_3)$.  We claim that $z\in A_2$. Suppose not, then $zx_{i}\in E(G)$ for some $i\notin \{1,2,3\}$. If $i=4$, then $\{x_4, z ,x_3, b_1, x_1\}$ induces a house, and if $i\neq 4$, then $\{ b_2, z,  x_3, x_4, x_i\}$ induces either a bull or a house. So $z\in A_2$ and this contradicts $z\notin B$. Therefore, $B$ is a connected homogeneous set.
\end{proof}

\end{document}
